% sections/difficultes.tex
% Phase 10: Difficultes rencontrees (RAPP-06)

\section{Difficult\'es rencontr\'ees}
\label{sec:difficultes}

Plusieurs probl\`emes techniques se sont pos\'es au cours du projet.

\paragraph{1. Complexit\'e algorithmique de la centralit\'e
d'interm\'ediarit\'e.}
Le calcul exact de la centralit\'e d'interm\'ediarit\'e est en
$O(|V| \cdot |E|)$, soit plusieurs heures pour notre graphe. On \'echantillonne
\`a la place $k = 500$ sommets pivots~\cite{freeman1977centrality} (graine 42),
ce qui ram\`ene le temps de calcul \`a quelques secondes tout en produisant un
classement fiable.

\paragraph{2. M\'etriques de chemins sur graphe d\'econnect\'e.}
Le diam\`etre et la longueur moyenne des chemins ne sont d\'efinis que sur
des graphes connexes. Le r\'eseau \'etant d\'econnect\'e (95 composantes
faiblement connexes), on calcule le diam\`etre sur la plus grande SCC et la
longueur moyenne sur la WCC par \'echantillonnage de 1\,000 paires.

\paragraph{3. Filtrage des transferts n\'egligeables (\textit{dust}).}
1\,598 ar\^etes (4\,\%) portent des volumes inf\'erieurs \`a $10^{-12}$~WETH
--- approbations automatis\'ees, pings de protocoles, erreurs d'arrondi. Sans
filtrage, elles gonflent artificiellement le nombre d'ar\^etes \`a volume nul
et biaisent le Gini. On les exclut de l'analyse distributionnelle (RQ6).

\paragraph{4. Conversion orient\'e/non-orient\'e pour la d\'etection de
communaut\'es.}
L'algorithme de Louvain~\cite{blondel2008louvain} requiert un graphe non
orient\'e, mais le r\'eseau de transferts est orient\'e. Appliquer Louvain
directement sur un \texttt{DiGraph} produit des r\'esultats incorrects. On
convertit donc le graphe en \texttt{Graph} non orient\'e en agr\'egeant les
poids (somme des volumes dans les deux sens) avant l'appel \`a
\texttt{louvain\_communities()}.

\paragraph{5. Encodage et typographie \LaTeX{} pour le fran\c{c}ais.}
La combinaison \texttt{cleveref} + \texttt{babel french} provoque des erreurs
de compilation (\textit{Missing} \texttt{\textbackslash endcsname}). On a
remplac\'e les commandes \texttt{\textbackslash cref\{\}} par la syntaxe
manuelle \texttt{section\~{}\textbackslash ref\{sec:...\}} dans l'ensemble
du rapport.
